\documentclass[a4paper,11pt]{article}
\usepackage{fancyhdr,graphics,graphicx,xy}
\usepackage{blkarray}
\usepackage{ifthen}
\usepackage{amssymb,amsmath,mathtools,color}
\pagestyle{fancy}
\usepackage{xparse}
\usepackage{nicematrix}


\textwidth = 6.5 in
\textheight = 9 in
\oddsidemargin = 0.0 in
\evensidemargin = 0.0 in
\topmargin = 0.0 in
\headheight = 0.0 in
\headsep = 0.0 in
\parskip = 0.2in
\parindent = 0.0in


%%%%%%%%%%%%%%%%%%%%%%%%%%%%%%%%%%%%%%%%%%
%%% General math macros
%%%%%%%%%%%%%%%%%%%%%%%%%%%%%%%%%%%%%%%%%%


\newcommand{\asstnumber}{2}
\newcommand{\assigntop}{
  \fancyhead[L]{}
  \setlength{\headheight}{37.40004pt}
  \fancyhead[C]{\Large Math 576 -- Fall 2023\\ Assignment \asstnumber}
  \fancyhead[R]{ }
}

\newtheorem{theorem}{Theorem}
\newtheorem{corollary}[theorem]{Corollary}
\newtheorem{definition}{Definition}

\begin{document}
\assigntop

\section{ Question 1}

If $X$ has the discrete topology, and it's compact, then 
$$
\bigcup_{x \in X} \{x\} \supseteq X
$$
is an open cover of $X$, since every singleton is an open set of $X$, then by compactness, we can take a finite subcover
$$ x_1 \cup x_2 \cup \ldots \cup x_n \supseteq X$$
thus $X$ is finite since it's cardinality is no more than $n$. Conversely, since $X$ is finite with the discrete topology, we can list out it's elements $\{x_1,x_2,\ldots, x_n\}$, then let $\{U_\alpha\}_{\alpha \in I}$ be an arbitrary open cover of $X$, then $U_{x} \subseteq \{U_\alpha \colon \alpha \in I\}$ be an open set that contains $x$, then 
$$ U_{x_1} \cup U_{x_2} \cup \ldots \cup U_{x_n}$$ 
is a finite subcover of $X$. 



\clearpage

\section{ Question 2 }
\subsection{(a)}
This is false, a straightforward example would be $X = Y = \mathbb{R}$ and $ X \sqcup Y = \mathbb{R} \sqcup \mathbb{R}$ which is Hausdorff, but the equivalence relation, defining $f \colon \mathbb{R}\setminus\{0\} \to \mathbb{R}\setminus\{0\}$ as $f(x) = x$ and $a \sim f(a)$ for all $a \in \mathbb{R}\setminus\{0\}$, we get a line with two origins, which is not hausdorff since a covering sequence towards the origin will have 2 limits. 
\subsection{(b)}
This is true, since let $x,y \in A$ such that $x \neq y$, then since, $x,y \in X$, there exists open sets $U,V$ open in $X$ such that $x \in U$ and $y \in V$ and $U \cap V = \emptyset$. But then $U \cap A$ and $V \cap A$ is open in $A$, and $x \in U \cap A$, $y \in V \cap A$ and $(U \cap A) \cap (V \cap A) = U \cap V = \emptyset$, thus $A$ is Hausdorff.

\clearpage

\section{problem 3}
Let $U \subseteq \mathbb{R}/ \sim$ open, then let $[x] \in U$, therefore, $x \in p^{-1}(U)$, where $p \colon x \mapsto [x]$ is the standard projection map. But since $p^{-1}(U)$ open in $\mathbb{R}$
\begin{equation}
	\exists \varepsilon > 0 \text{ such that } B_\varepsilon(x) (x-\varepsilon,x+\varepsilon) \subseteq U
\end{equation}
Then let $y \in \mathbb{R}$, since the rationals are dense in the reals, $\exists q \in \mathbb{Q}$ such that 
\begin{equation}
	|(x-y)-q| = |x-(y+q)| < \varepsilon
\end{equation}
But then 
\begin{equation}
	(y+q) \in B_\varepsilon(x) \subseteq p^{-1}(U) \implies [y+q] \in U
\end{equation}
but observe that
\begin{equation}
	(y+q) - y = q \in \mathbb{Q} \implies y+q \sim y \implies [y] \in U
\end{equation}
but then 
\begin{equation}
	y \in p^{-1}(U)
\end{equation}
Since $y$ is an arbitrary real number, we have that 
\begin{equation}
	p^{-1}(U) = \mathbb{R} \implies U = \mathbb{R}/\sim \text{ since $p$ surjective}
\end{equation}
That means the only open sets in $\mathbb{R}/\sim$ are $\mathbb{R}/\sim$ and $\emptyset$, thus $\mathbb{R}/\sim$ has the indiscrete topology. Intuitively, this means that we have squished the reals number line by quite a a lot. 
\clearpage

\section{question 4}
\subsection{(a)}
We verify the axioms for a basis,
\begin{itemize}
	\item[(i)] let $x \in S^\mathbb{N}$, then let $p = \{x_1,\ldots,x_N\}$ for any $N \geq 1$, then $x \in U_p$. 
	\item[(ii)] Let $U_x, U_y \in \mathcal{B}$, then let $p \in U_x \cap U_y$. Noting that $U_x \cap U_y$ contains sequences of $S^\mathbb{N}$ that agree with $x$ \textbf{and} $y$ up to $n$
		\begin{equation}
			N = \max_{i} (x_i = y_i)
		\end{equation}
		then define 
		\begin{equation}
			q = \{x_1,\ldots,x_N\}
		\end{equation}
		We can claim that 
		\begin{equation}
			U_q \subseteq U_x \cap U_y
		\end{equation}
		since any $s \in U_q$ agrees with $q$ up to $N$, but since $q_i=x_i=y_i$ for $1 \leq i \leq N$, then $s \in U_x$ and $s \in U_y$, and clearly, $p \in U_q$. 
		
\end{itemize}
\subsection{(b)}
let $U_p$ be any open set containing $x$ with $p$ as a representative, that means $x$ agrees with $p$ up to some $N$, then clearly $\forall n \geq N$, $y_n \in U_p$ since taking any $y_k$ with $k \geq N$, then $x_i=y_i=p_i$ for $i \leq N$. Thus $y_k$ agrees with $p$ up to $N$ as well. Since $U_p$ is arbitrary $y_n$ converges to $x$. 
\subsection{(c)}
Let $x,y \in S^\mathbb{N}$, then let 
\begin{equation}
	N = \min \{i \in \mathbb{N} \colon x_i \neq y_i\} 
\end{equation}
then let $p = \{x_1,\ldots,x_{N}\}$ and $q = \{y_1,\ldots,x_{N}\}$, then clearly $x \in U_p$ and $y \in U_q$. $U_p \cap U_q = \emptyset$, because any sequence $p^\prime$ in $U_p$ agrees with $x$ up to $x_{N}$, but $p^\prime_{N} \neq y_{N}$ and vice-versa, therefore, $X$ is Hausdorff. 
\subsection{(d)}
We want to show that any set with at least two points is disconnected. Let $A \subseteq X$ have at least 2 elements $x,y$ with $x \neq y$, but then letting 
\begin{equation}
	N = \min \{ i \colon x_i \neq x_y\}
\end{equation}
This is the first index at which $x$ and $y$ are not the same, then let
\begin{align}
	U = \{p \in S^\mathbb{N} \colon p_N \neq x_N\}\\
	V = \{p \in S^\mathbb{N} \colon p_N = x_N\}
\end{align}
These are open sets, since letting any $q \in U$, then $\{ p \in S^{\mathbb{N}} \colon p_i = q_i \text{ up to $N$} \} \subseteq U$ is an open set containing $q$. $U$ and $V$ are clearly disjoint, and their union partition $X$.Since $U$ and $V$ partition the whole space, they partition $A$. Furthermore, $U \cap A$ and $V \cap A$ are disjoint and non empty since $y \in U$ and $x \in V$, thus $A$ is disconnected. Thus taking the contrapositive, any connected set is a singleton. 
	
\clearpage

\section{question 5}

\subsection{(a)}
Let $Z$ be a topological space with continuous maps $i \colon X \to Z$ and $j \colon Y \to Z$ and $A \subseteq X$, and $B \subseteq Y$, and a map $f \colon A \to B$, such that $i(a) = j(f(a))$ with an equivalence relation $a \sim f(a)$ then there exists a unique continuous map $\bar{h} \colon X \sqcup Y/\sim \to Z$ such that the SQUARE commutes. 

	Since $i$ and $j$ are continuous, then by universal property of disjoint union, there exists unique continuous map $h \colon X \sqcup Y \to Z$. We can easily verify that $h$ is continuous by setting $h_{|X} = i(x)$ and $h_{|Y} j(x)$, then taking an open set $U \subseteq Z$, since $Im(i) \cup Im(j) = X \sqcup Y$,  we get $h^{-1}(U) = i^{-1}(U) \cup j^{-1}(U)$ is open in $X$. Then let $a \sim b$ in $X$, then $a \sim f(a)$, then we have that $h(a) = i(a) = j(f(a)) = h(f(a)) = h(b)$, thus we can use the universal property of quotient spaces to get a unique continuous map $\bar{h} \colon X \sqcup Y /\sim \to Z$ such that the SQUARE commutes. 
\subsection{(b)}
Let $f \colon A \to B$ such that $f \colon 1 \mapsto 5$ and the identity elsewhere, and define the equivalence relation of $a \sim f(a)$ on $X \sqcup Y$. We try to show that $X \sqcup Y /\sim $ is homeomorphic to $[0,2]$. Define $i \colon [0,1] \to [0,1]$ and $j \colon [5,6] \to [1,2]$ 
\begin{align}
	i(x) &= x \\
	j(x) &= x-4 
\end{align}
Clearly here we have that $i(1) = 1 = j(5) = j(f(1))$. Then we can define a unique continuous extension onto $h \colon [0,1] \sqcup [5,6] \to [0,2]$ just like in (a), then since $h(1) = i(1) = j(5) = j(f(1)) = h(f(1))$. $h$ here is a quotient map that agrees with the glueing which produces a homeomorphism with the map $\bar{h} \colon X \sqcup Y / \sim \to [0,2]$ 
\subsection{(c)}
Define a quotient map $h \colon X \sqcup Y = D^2 \times \{1\} \cup D^2 \times \{0\} \to S^2$ as
\begin{align}
	(x_1,y_1) \mapsto (x_1,y_1,\sqrt{1-|x|^2})\\
	(x_2,y_2) \mapsto (x_2,y_2,-\sqrt{1-|x|^2})
\end{align}
This is continuous and has a continuous inverse, and it is surjective. It also satisfies $x \sim f(x) \implies h(x) = h(f(x))$ since we map every point on the boundary of the discs to the equator of the sphere, therefore by (a), we have $ X \sqcup Y/\sim$ homeomorphic to $S^2$.




\end{document}
